% figures/fig2-ir-expr.tex — IR expressiveness comparison
\begin{figure}[t]
\centering
\begin{tikzpicture}[
  >=Stealth,
  scale=0.75, every node/.style={scale=0.75},
  panel/.style={draw, rounded corners=2pt, inner sep=4pt, align=left,
                minimum width=4.6cm, minimum height=2.6cm},
  lbl/.style={font=\footnotesize\bfseries, anchor=south, inner sep=2pt},
  code/.style={font=\scriptsize\ttfamily, align=left, anchor=north west},
  note/.style={font=\scriptsize, align=center, text width=4.4cm},
]
  % --- Left: LLVM IR ---
  \node[panel] (left) {};
  \node[lbl] at (left.north) {General-purpose IR (LLVM)};

  \node[code] at ([shift={(4pt,-4pt)}]left.north west) {%
    \%r = call \{i256, i1\}\\
    \quad @llvm.uadd.with.overflow\\
    \quad\quad .i256(\%x, \%y)\\[1pt]
    \%sum\enspace = extractvalue \%r, 0\\
    \%cout = extractvalue \%r, 1%
  };

  \node[note, anchor=north] at ([yshift=-3pt]left.south) {%
    Carry bit carried indirectly via\\a multi-return tuple};

  % --- Right: dMIR (separated to avoid note overlap) ---
  \node[panel, right=1.6cm of left] (right) {};
  \node[lbl] at (right.north) {dMIR};

  \node[code] at ([shift={(4pt,-4pt)}]right.north west) (rcode) {%
    (sum, cout) = ADC(x, y, cin)%
  };

  \node[font=\scriptsize, align=left, anchor=north west]
    at ([yshift=-2pt]rcode.south west) {%
    $: \mathrm{BV}_{256}\!\times\!\mathrm{BV}_{256}\!\times\!\mathrm{BV}_{1}$\\
    $\;\;\to \mathrm{BV}_{256}\!\times\!\mathrm{BV}_{1}$%
  };

  \node[note, anchor=north] at ([yshift=-3pt]right.south) {%
    Carry bit expressed directly\\as a separate bit-vector};

\end{tikzpicture}
\caption{Expressiveness comparison of general-purpose IR and dMIR on
carry-propagating addition.}
\label{fig:ir-expr}
\end{figure}
